%%% Document Formatting
\documentclass[11pt,fleqn]{article}
\usepackage[a4paper,
            bindingoffset=0.2in,
            left=0.75in,
            right=0.75in,
            top=0.8in,
            bottom=0.8in,
            footskip=.25in]{geometry}
\setlength\parindent{10pt}

%%% Imports
\usepackage{amsmath, amsfonts, amssymb, amsthm}
\usepackage{dsfont}
\numberwithin{equation}{section}
\usepackage{mathtools}
\usepackage{cancel}
\usepackage{empheq}
\usepackage{algorithm}
\usepackage{algpseudocode}  % preamble

% Visualization
\usepackage{graphicx}


%%% Formatting
\usepackage{hyperref}
\hypersetup{
    colorlinks=true,
    linkcolor=blue,
    filecolor=magenta,      
    urlcolor=cyan,
    pdftitle={SCSI Lecture Notes},
    pdfpagemode=FullScreen,
}
\urlstyle{same}

%%% Custom Environments
 

% theorems
% Dependencies
\usepackage[usenames,dvipsnames]{xcolor}
\usepackage{amsmath, amsfonts, mathtools, amsthm, amssymb}
\usepackage{thmtools}
\usepackage[framemethod=TikZ]{mdframed}
\mdfsetup{skipabove=1em,skipbelow=0em}

\theoremstyle{definition}


\makeatother
\usepackage{thmtools}
\usepackage[framemethod=TikZ]{mdframed}
\mdfsetup{skipabove=1em,skipbelow=0em}


\theoremstyle{definition}

\declaretheoremstyle[
    headfont=\bfseries\color{ForestGreen!70!black}, bodyfont=\normalfont,
    mdframed={
        linewidth=2pt,
        rightline=false, topline=false, bottomline=false,
        linecolor=ForestGreen, backgroundcolor=ForestGreen!5,
    }
]{thmgreenbox}

\declaretheoremstyle[
    headfont=\bfseries\color{NavyBlue!70!black}, bodyfont=\normalfont,
    mdframed={
        linewidth=2pt,
        rightline=false, topline=false, bottomline=false,
        linecolor=NavyBlue, backgroundcolor=NavyBlue!5,
    }
]{thmbluebox}

\declaretheoremstyle[
    headfont=\bfseries\color{NavyBlue!70!black}, bodyfont=\normalfont,
    mdframed={
        linewidth=2pt,
        rightline=false, topline=false, bottomline=false,
        linecolor=NavyBlue
    }
]{thmblueline}

\declaretheoremstyle[
    headfont=\bfseries\color{RawSienna!70!black}, bodyfont=\normalfont,
    mdframed={
        linewidth=2pt,
        rightline=false, topline=false, bottomline=false,
        linecolor=RawSienna, backgroundcolor=RawSienna!5,
    }
]{thmredbox}

\declaretheoremstyle[
    headfont=\bfseries\color{RawSienna!70!black}, bodyfont=\normalfont,
    numbered=no,
    mdframed={
        linewidth=2pt,
        rightline=false, topline=false, bottomline=false,
        linecolor=RawSienna, backgroundcolor=RawSienna!1,
    },
    qed=\qedsymbol
]{thmproofbox}

\declaretheoremstyle[
    headfont=\bfseries\color{NavyBlue!70!black}, bodyfont=\normalfont,
    numbered=no,
    mdframed={
        linewidth=2pt,
        rightline=false, topline=false, bottomline=false,
        linecolor=NavyBlue, backgroundcolor=NavyBlue!1,
    },
]{thmexplanationbox}

\declaretheoremstyle[
    headfont=\bfseries\color{Plum!70!black}, bodyfont=\normalfont,
    mdframed={
        linewidth=2pt,
        rightline=false, topline=false, bottomline=false,
        linecolor=Plum, backgroundcolor=Plum!5,
    }
]{thmpurplebox}

\declaretheoremstyle[
    headfont=\bfseries\color{Sepia!70!black}, bodyfont=\normalfont,
    mdframed={
        linewidth=2pt,
        rightline=false, topline=false, bottomline=false,
        linecolor=Sepia, backgroundcolor=Sepia!5,
    }
]{thmsepiabox}

\declaretheoremstyle[
    headfont=\bfseries\color{ForestGreen!70!black}, bodyfont=\normalfont,
    mdframed={
        linewidth=2pt,
        rightline=false, topline=false, bottomline=false,
        linecolor=ForestGreen
    }
]{thmgreenline}

% \declaretheoremstyle[headfont=\bfseries, bodyfont=\normalfont, mdframed={ nobreak } ]{thmgreenbox}
% \declaretheoremstyle[headfont=\bfseries, bodyfont=\normalfont, mdframed={ nobreak } ]{thmredbox}
% \declaretheoremstyle[headfont=\bfseries, bodyfont=\normalfont]{thmbluebox}
% \declaretheoremstyle[headfont=\bfseries, bodyfont=\normalfont]{thmblueline}
% \declaretheoremstyle[headfont=\bfseries, bodyfont=\normalfont, numbered=no, mdframed={ rightline=false, topline=false, bottomline=false, }, qed=\qedsymbol ]{thmproofbox}
% \declaretheoremstyle[headfont=\bfseries, bodyfont=\normalfont, numbered=no, mdframed={ nobreak, rightline=false, topline=false, bottomline=false } ]{thmexplanationbox}

\declaretheorem[style=thmgreenbox, name=Definition]{definition}
\declaretheorem[style=thmgreenbox, name=Assumption]{assumption}
\declaretheorem[style=thmbluebox, numbered=no, name=Example]{eg}
\declaretheorem[style=thmredbox, name=Proposition]{prop}
\declaretheorem[style=thmredbox, name=Theorem]{theorem}
\declaretheorem[style=thmredbox, name=Lemma]{lemma}
\declaretheorem[style=thmredbox, numbered=no, name=Corollary]{corollary}
\declaretheorem[style=thmpurplebox, name=Problem]{problem}
\declaretheorem[style=thmsepiabox, numbered=no, name=Fact]{fact}
\declaretheorem[style=thmgreenline, name=Property]{property}

\declaretheorem[style=thmproofbox, name=Proof]{replacementproof}
\renewenvironment{proof}[1][\proofname]{\vspace{-10pt}\begin{replacementproof}}{\end{replacementproof}}


\declaretheorem[style=thmexplanationbox, name=Proof]{tmpexplanation}
\newenvironment{explanation}[1][]{\vspace{-10pt}\begin{tmpexplanation}}{\end{tmpexplanation}}

\declaretheorem[style=thmblueline, numbered=no, name=Remark]{remark}
\declaretheorem[style=thmblueline, numbered=no, name=Note]{note}

\newmdenv[ 
  topline=false,
  bottomline=false,
  skipabove=\topsep,
  skipbelow=\topsep
]{sidework}

\title{Mathematical Modeling with Data \\[0.3em]
       \large NYU MATH-GA 2830-003 --- Andrew Stuart --- Fall 2026}
\author{Clark Miyamoto}
\date{\today}

\begin{document}

\maketitle

We follow a modified version of \url{https://arxiv.org/abs/2410.10523}. The purpose of this course is solve inverse problems (and a subclass of said problems, data assimilation). Our schedule is:
\begin{enumerate}
	\item Lectures 1-6: Formalism \& Recap of Machine Learning.
	\item Lectures 7-10: Theory of inverse problems / problem statement.
	\item Then technqiues on sovling inverse problems / data assimilation using ML / classical techniques.
	\item Stating LLMs \& Contrastive Learning in this framework.
\end{enumerate}
Prof. Stuart will present/emphasize the material slightly different than the lecture notes.

\tableofcontents

\newpage

\section{Introduction}

\subsection{Scoring Rules}
Let $\rho,\rho' \in \mathcal P(\mathbb R^d)$ and $u , v, w \in \mathbb R^d$. To simplify our lives, we'll mainly work with measure which emit probability density functions (pdfs) $\left\{ \rho \in L^1(\mathbb R^d; [0,\infty)) : \int \rho = 1 \right\} \subset \mathcal P(\mathbb R^d)$.

\paragraph{Metric Spaces for Probability Distributions} 
We'll be focusing on ways to assess distances between probability distributions. We'll go over (1) KL-divergence, important in theory \& algorithmically (2) Hellinger metric (3) Energy distance.

\begin{definition}[Negative Log Likelihood]
	Consider the function $\textsf{LS}: \mathcal P(\mathbb R^d) \times \mathbb R^d \to [-\infty, \infty]$,
\begin{align}
	\textsf{LS}(\rho, v) = - \log \rho(v)
\end{align}
The notation is because this will becomes a \emph{\textsf{S}coring rule}.
\end{definition}
Intuitively, the smaller this is, the more likely $v \sim \rho$, and vice versa. However note, in high-dimensions this might not be the case (think typical set).

\begin{definition}
	[KL Divergence] Consider the function $\textsf{KL}: \mathcal P(\mathbb R^d) \times \mathcal P(\mathbb R^d) \to [0, \infty]$
	\begin{align}
		\textsf{KL}(\rho\ \| \rho) = \mathbb E_{u \sim \rho'} \left[ \log \frac{\rho'(u)}{\rho(u)} \right] = \mathbb E_{u \sim \rho'} \textsf{LS}(\rho, u) - \mathbb E_{u \sim \rho'} \textsf{LS}(\rho', u)
	\end{align}
\end{definition}
Properties
\begin{enumerate}
	\item (Non-negative) $\textsf{KL} \geq 0$.
	\item (Identifiable (?)) $\textsf{KL}(\rho \| \rho') = 0 \iff \rho = \rho'$.
\end{enumerate}
\begin{proof} 
	\begin{align}
		\textsf{KL}(\rho' \| \rho) & = - \mathbb E_{\rho'} \log \frac{\rho}{\rho'}\\
		& \geq - \log \mathbb E_{\rho'} \frac{\rho}{\rho'} & \text{Jensen}\\
		& = - \log 1 = 0
	\end{align}
	Further more, because $\log x$ strictly convex, you only get equality i.f.f. $\rho'(u)/\rho(u) = 1$ (almost everywhere).
\end{proof}

\begin{definition}
	[Scoring Rule] Let $s: \mathcal P(\mathbb R^d) \times \mathbb R^d \to \mathbb R$ is a \textbf{scoring rule}, intuitively it will determine the "quality" of sample drawn from an associated measure. The \textbf{expected scoring rule} is a function $\bar s : \mathcal P(\mathbb R^d) \times \mathcal P(\mathbb R^d) \to \mathbb R^d$ given by $\bar s(\rho, \rho') = \mathbb E_{u \sim \rho'}s(\rho, u)$.
	\\
	\\
	The notation $\bar s$ comes from averaging (bar operator) over the argument $s(\rho, \cdot)$.
\end{definition}


So for example, an expected scoring rule is the KL-divergence. Expected scoring rules are nice in machine learning because you can estimate $\mathbb E_{u \sim \rho'}$ numerically.

\begin{definition}
	[Proper Scoring Rule] An expected scoring rule is \textbf{proper} if
	\begin{align}
		\bar s(\rho', \rho') \leq \bar s(\rho, \rho'), \quad \forall \rho,\rho' \in \mathcal P(\mathbb R^d)
	\end{align}
	It is \textbf{strictly proper} if $\bar s(\rho', \rho') = \bar s(\rho, \rho') \iff \rho = \rho'$
\end{definition}


Let $\rho \in \mathcal P(\mathbb R)$ (yes, 1 dimension).
\begin{definition}
	[Cumulative Distribution Function (CDF)] 
	\begin{align}
		F_\rho(u) := \int_{-\infty}^u \rho(w) \, dw
	\end{align}
\end{definition}
\begin{eg}
	Let $\rho = \textsf{Unif}[0,1]$. Then $F_\rho(u) = \begin{cases}
		0 & u <0\\
		u& u \in [0, 1)\\
		1 & u \geq 1
	\end{cases}$
\end{eg}
\begin{eg}
	Let $\rho = \delta_v$, then $F_\rho(u) = \begin{cases}
		0& u < v\\
		1 & u \geq v
	\end{cases} = \mathds 1_{u \geq v}$. Where $\mathds 1$ is the indicator function.
\end{eg}

\begin{definition}
	[Continuous Rank Probability Score] Let $\rho \in \mathcal P(\mathbb R)$ and $v \in \mathbb R$.
	\begin{align}
		\textsf{CRPS}(\rho, v) = \| F_\rho(\cdot) - \mathds 1_{\cdot \geq v}\|^2_{L^2(\mathbb R)}
	\end{align}
	This has been used in weather forecasting, where $v$ is measurements of pressure, velocity, etc. at a single weather station.
\end{definition}
Notice this is indeed a scoring rule.
\begin{definition}
	[Quantile Score] Let $q_\rho: (0,1) \to \mathbb R$, defined by $F_\rho(q_\rho(\alpha)) = \alpha$ ($\forall \alpha \in (0,1)$). 
\end{definition}
So $q_\rho$ is essentially the inverse of the CDF $F_\rho$; note, the CDF isn't always invertible. For example, when the CDF is flat or has infinite slope, its inverse is not unique. To break this degeneracy, we'll take the inverse which is intuitive via flipping \& rotation the function graphically.
\begin{eg} 	\begin{enumerate}
		\item $\rho = \textsf{Unif}[0,1]$, then $q_\rho(\alpha) = \alpha$
		\item $\rho = \delta_v$, then $q_\rho(\alpha) = v$
	\end{enumerate}
	Notice $\alpha \in (0,1)$.
\end{eg}

\begin{definition}
	Let
	\begin{align}
		h_\alpha(u) = \begin{cases}
			(1-\alpha) |u|, & u < 0\\
			\alpha |u|, & u \geq 0
		\end{cases}
	\end{align}
\end{definition}
\begin{lemma}
	Let $L_{\alpha, \rho}(\theta) = \mathbb E_{v \sim \rho} [h_\alpha(v - \theta)]$. Then $q_\rho(\alpha) = \arg \min_{\theta \in \mathbb R} L_{\alpha, \rho}(\theta)$. Where $q_\rho$ is the quantile.
\end{lemma}
\begin{definition}
	[Quantile Score] A quantile score at level $\alpha$ if
	\begin{align}
		\textsf{QS}_\alpha(\rho, v) = h_\alpha(v - q_\rho(\alpha))
	\end{align}
	The expected quantile score is
	\begin{align}
		\bar{\textsf{QS}}_\alpha(\rho, \rho') = \mathbb E_{v \sim \rho'} h_\alpha(v - q_\rho(\alpha))
	\end{align}
\end{definition}
\begin{lemma}
	$\textsf{QS}_\alpha$ is a proper scoring rule.
\end{lemma}

\begin{theorem}
	\begin{align}
		\textsf{CRPS}(\rho, v) = 2 \int_0^1 \textsf{QS}_\alpha(\rho, v) \, d\alpha
	\end{align}
\end{theorem}

\subsection{Divergence}
\begin{definition}
	[Divergence] Let $D : \mathcal P(\mathbb R^d) \times \mathcal P(\mathbb R^d) \to \mathbb R$ denote a \emph{(statistical) divergence}, if 
	\begin{itemize}
		\item $D(\rho_1 \| \rho_2) \geq 0$ 
		\item $D(\rho_1 \| \rho_2) = 0 \iff \rho_1 = \rho_2$
	\end{itemize}
	In contrast to a metric, it is not symmetric in its arguments. Notice all distances are also divergences.
\end{definition}
\begin{theorem}
	If $s$ is a strictly proper scoring rule, then $D_s(\rho_1, \rho_2)$ is a divergence. In particular
	\begin{align}
		D_s(\rho_1 \| \rho_2) = \bar s(\rho, \rho') - \bar s(\rho' , \rho')
	\end{align}
\end{theorem}
\begin{eg}
	Let $s(\rho, v) = - \log \rho(v)$ which is strictly proper. This corresponds to divergence $D_{KL} (\cdot \| \cdot)$.
\end{eg}

\begin{definition}
	Let $f: [0, \infty) \to (-\infty, \infty]$ be a convex function, and $f(0) = \lim_{t \to 0^+} f(t)$. Then the $f$-divergence is 
	\begin{align}
		D_f(\rho_1 \| \rho_2) = \mathbb E_{u \sim \rho_2} f \left( \frac{\rho_1(u)}{\rho_2(u)} \right)
	\end{align}
\end{definition}
\begin{theorem}
	Let $f$ be a strictly convex function, then $D_f(\cdot \| \cdot)$ is indeed a divergence.
\end{theorem}
\begin{proof}
	Showing $\geq 0$
	\begin{align}
		D_f(\rho_1 \| \rho_2) & = \mathbb E_{\rho_2} \left[ f \left( \frac{\rho_1(u)}{\rho_2(u)} \right) \right]\\ 
		& \geq f \left( \mathbb E_{\rho_2} \left[ \frac{\rho_1}{\rho_2} \right] \right) & \text{Jensen}\\
		& = f(0) & \rho_2 \text{ cancel} \\
		& = 0 & \text{Definition of } f
	\end{align}
	Now showing $\rho_1 = \rho_2 \iff D_f = 0$. Proving $\implies$ is trivial. Now doing the other way
\end{proof}
Note, $\textsf{TV}$ distance is also a divergence, but the choice of $f(t) = \frac{1}{2} | t-  1|$ is not strictly convex. Hence why he defined $f$-divergences to just be convex functions.

\begin{eg}
	$f(t) = t \log t$ corresponds to the $\textsf{KL}$ divergence. You can also get the reverse $\textsf{KL}$ by using $f(t) = - \log t$. \\
	\\
	Another divergence is $f(t) = (t-1)^2$ which corresponds to $\chi^2$.
\end{eg}

\begin{definition}
	[Metric] A metric $D$ is a divergence which fits two additional criteria. Let $\rho_1, \rho_2, \rho_3 \in \mathcal P(\mathbb R^d)$.
	\begin{enumerate}
		\item Triangle inequality: $D(\rho_1, \rho_3) \leq  D(\rho_1 , \rho_2) + D(\rho_2, \rho_3)$
		\item Symmetric: $D(\rho_1, \rho_2) = D(\rho_2, \rho_2)$
	\end{enumerate}
\end{definition}
Why not define these using norms? Well probability densities don't form a vector space.

\begin{definition}
	[Energy Score] For $\beta \in (0, 2]$
	\begin{align}
		\textsf{ES}_\beta(\rho, v) = \mathbb E_{u \sim \rho} |u - v|^\beta - \frac{1}{2} \mathbb E_{(a, a') \sim \rho \otimes \rho} |a -a'|^\beta
	\end{align}
	where $(a, a') \sim \rho \otimes \rho$ means $a,a'$ are drawn iid from the same distribution $\rho$. The divergence corresponding to the energy score is given by
	\begin{align}
		D^2_{E, \beta}(\rho_1, \rho_2) : = \bar{\textsf{ES}}_{\beta}(\rho, \rho') -  \bar{\textsf{ES}}_{\beta}(\rho', \rho') = \mathbb E_{u,v} |u - v|^\beta - \frac{1}{2} \mathbb E_{u,u'} |u- u'|^\beta - \frac{1}{2} \mathbb E_{v,v'} |v - v'
	|^\beta
	\end{align}
	where all expectations are taken over $\rho \otimes \rho$. 
\end{definition}
In this class, we'll typically take $\beta = 1$.
Interestingly this divergence gets sweetened into a norm (over characteristic functions).
\begin{theorem}
	Let $\rho_1, \rho_2 \in \mathcal P(\mathbb R^d)$ and $\beta \in (0,2)$. Then $\exists \,  C = C(\beta, d)$  (this constant blows up as $\beta \to 2$)
	\begin{align}
		\frac{1}{2} D^2_{E, \beta}(\rho_1, \rho_2) = C \int_{\mathbb R^d} \frac{| \varphi_\rho(u) - \varphi_{\rho'}(u)|^2}{|u|^{d + \beta}} du
	\end{align}
	where $\varphi_\rho$ is the characteristic function of $\rho$. This is an L2 norm over characteristic functions.  Therefore $\frac{1}{2}D^2_{E, \beta}$ is a divergence, and $D_{E, \beta}$ is a metric.
\end{theorem}
\begin{proof}
	Proof is magical. Andrew doesn't go over it in class.
\end{proof}

\begin{sidework}
	Notation: $\textsf{ES}(\rho, v)$ = $\textsf{ES}_1(\rho, v)$. \\
	\\
	Recall: $\textsf{CPRS}(\rho, v) = \int |F_\rho(z) - \mathds 1_{z \geq v}(z) |^2 \, dz$, where $F_\rho$ is the cumulative function corresponding to $\rho$. This measures how far your deviate from $\rho = \delta_v$.
\end{sidework}

\begin{theorem}
	In $d = 1$, $\textsf{CRPS} = \textsf{ES}$.
\end{theorem}
\begin{proof}
	 First note
	 \begin{align}
	 	|u-v| = \int_v^\infty \mathds 1_{z \leq u} \, dz + \int_{-\infty}^v \mathds 1_{z > u} \, dz
	 \end{align}
	 By Fubini
	 \begin{align}
	 	\mathbb E_{u \sim \rho}|u-v| & =  \int \rho(u) |u-v|\ , du\\
	 	& = \int \rho(u) \left( \int_v^\infty \mathds 1_{z \leq u} \, dz + \int_{-\infty}^v \mathds 1_{z > u} \, dz \right)\\
	 	& = \int_v^\infty [1 - F_\rho(z)] \, dz + \int_{-\infty}^u F_\rho(z) \, dz\\
	 	& = \int_{\mathbb R} [1-F_\rho(z)] \mathds 1_{z \geq v} + F_\rho(z) \underbrace{\mathds 1_{z < v}}_{1 - \mathds 1_{z \geq v}} \, dz
	 \end{align}
	 Ok, now we can integrate
	 \begin{align}
	 	\frac{1}{2}\mathbb E_{(u,v) \sim \rho^{\otimes 2}} |u-v| = \int [1-F_\rho(z)] \, F_\rho(z) \, dz
	 \end{align}
	 Putting it all together
	 \begin{align}
	 	\textsf{ES}(\rho, v) & = \mathbb E_{u \sim \rho} |u -v| - \frac{1}{2} \mathbb E_{u,u' \sim \rho^{\otimes 2}}  |u - u'|\\
	 	& = \int_{\mathbb R} [F_\rho(z)^2 - 2 F_\rho(z) \mathds 1_{z \geq v} + \mathds 1_{z \geq v}^2] \, dz\\
	 	& = \int_{\mathbb R} [F_\rho(z) - \mathds 1_{z \geq v}]^2 \,dz = \textsf{CRPS}(\rho, v)
	 \end{align}
\end{proof}

\begin{definition}
	Let $\rho, \rho' \in \mathcal P(\mathbb R^d)$ be probability density functions. For $1 \leq p < \infty$, define the metric
	\begin{align}
		D_p(\rho, \rho')^p & = \frac{1}{2^{1/p}} \int |\rho(u)^{1/p} - \rho'(u)^{1/p}|^p \, du\\
		 & = \| |\rho|^{1/p} - |\rho'|^{1/p}\|_{L^p}^p
	\end{align}
\end{definition}
\begin{eg}
	$p=1$ is TV distance, and $p=2$ is the Hellinger Metric.
\end{eg}

\begin{theorem}
	For $p \in (1, \infty)$, there exists $K = K(p)$, such that
	\begin{align}
		\frac{1}{K}D_1(\rho,\rho') \leq D_p (\rho, \rho') \leq D_1(\rho, \rho')^{1/p}
	\end{align}
\end{theorem}
\begin{proof}
	Homework.
\end{proof}

\begin{eg}
	If $\textsf{TV} = \mathcal O(\epsilon)$ (which is $p=1$), then this implies $D_{Hell} = \mathcal O(\epsilon^{1/2})$.\\
	\\
	If $D_{Hell} = \mathcal O(\epsilon)$, then this implies $D_{TV}= \mathcal O(\epsilon)$. Meaning the Hellinger metric is much stronger than the TV distance.
\end{eg}

\newpage
\section{Transport}
 
\begin{definition}
	[Pushforward] Let $p \in \mathcal P(\mathbb R^d)$ and $p' \in \mathcal P(\mathbb R^{d'})$, and $T : \mathbb R^d \to \mathbb R^{d'}$. The notation $p' = T_\# p$   (called the pushforward) means
	\begin{itemize}
		\item If $x \sim p$ then $T(x) \sim p'$
		\item $\forall$ bounded \& measurable $\phi: \mathbb R^d \to \mathbb R$
		\begin{align}
			\int \phi(y) p'(y) \, dy = \int \phi(y) (T_\# p)(y) \, dy = \int \phi(T(x)) p(x) \,dx
		\end{align}
	\end{itemize}
\end{definition}

\paragraph{Transport Metrics}
\begin{definition}
	[Coupling] Let $p \in \mathcal P(\mathbb R^d)$ and $p' \in \mathcal P(\mathbb R^{d'})$. Let $\pi \in \mathcal P(\mathbb R^d \times \mathbb R^d)$ is a \textbf{coupling} of $p, p'$ if it has the specified marginals
	\begin{align}
		\int \pi(u,v) \, dv = p(u), \quad \int \pi(u,v) \, du = p'(v)
	\end{align}
	We'll denote $\Pi_{p,p'} = \left\{ \text{set of all couplings of } p,p' \right\}$.
\end{definition}

\begin{definition}
	[Kantorovich Formulation of Optimal Transport] Specify a cost function $c: \mathbb R^d \times \mathbb R^d \to \mathbb R_{\geq 0}$.
	\begin{align}
		\left\{ \pi^* \in \arg \min_{\pi \in \Pi_{p,p'}} \int_{\mathbb R^d \times \mathbb R^d} c(u,v) \pi(u,v) \, du \, dv \right\}
	\end{align}
\end{definition}
\begin{eg}
	Take $c(u,v) = |u-v|^2$
\end{eg}
\begin{remark}
	When you're trying to find the optimal transport between two atomic measures, this is a linear programming problem.\\
	\\
	To see this, map $c \mapsto \textsf{c}, \pi \mapsto \textsf{P}$ and $p,p' \mapsto \textsf{p,p'}$. You can now see the Kantorovich problem as
	\begin{align}
		&\inf_{\textsf{P}} \langle \textsf{c},\textsf{P}\rangle\\
		& \text{subject to } \textsf{P} \mathds 1 = \textsf{p}, \text{ and } \mathds 1^T \textsf{P} = \textsf{p'}
	\end{align}
	So $\mathds 1$ is equivalent to the marginalization operator in function space.
\end{remark}

\begin{definition} [Wasserstein-2 Distance/Metric]
	\begin{align} 
		W_2(p,p')^2 = \inf_{\pi \in \Pi_{p,p'}} \int |u-v|^2 \pi(u,v) \, du \, dv
	\end{align}
\end{definition}
This is a true metric (always $\geq 0$, $=0$ if and only if arguments are equal, symmetric, and satisfies triangle inequality).


\begin{definition}
	[Monge Formulation of Optimal Transport] Let $g : \mathbb R^d \to \mathbb R^d$. The Monge problme is to find a coupling $\pi \in \Pi_{p,p}$ s.t. $\pi = (\textsf{id}, g)_\# p$. We'll denote the solution space as $\Pi^{\textsf{Monge}}_{p,p'}$.
	\\
	\\
	The notation $(\textsf{id}, g): \mathbb R^d \to \mathbb R^d \times \mathbb R^d$.
\end{definition}

\begin{eg}
	Consider $y = g(x) + \eta$ where $x \sim p, \eta \sim \mathcal N(0, \Sigma)$ (s.t. $x \perp \eta$). Then
	\begin{align}
		y \sim p' = (g_\# p) * \mathcal N(0, \Sigma)
	\end{align}
\end{eg}
\begin{lemma}
	If $\pi \in \Pi_{p,p'}^{\textsf{Monge}}$, then 
	\begin{enumerate}
		\item $\pi(u,v) = p(u) \times \delta_{g(u)}(v)$
		\item $p' = g_\# p$.
	\end{enumerate}
\end{lemma}
\begin{proof}
	Let $\Omega = \mathbb R^d \times \mathbb R^d$. 
	\begin{enumerate}
		\item \emph{We want to show} $\forall$ bounded \& measurable $\psi : \Omega \to \mathbb R$
	\begin{align}
		\int \psi(u,v) p(u) \delta_{g(u)}(v) du\, dv = \int \psi(u,v) \left( (\textsf{id}, g)_\# p \right) (u,v) \, du \, dv
	\end{align}
	\emph{Starting proof} from the LHS
	\begin{align}
		\textsf{LHS} & = \int \left( \int \psi(u,v) \delta_{g(u)}(v) \, dv \right) p(u) \, du & \text{Boundedness}\\
		& = \int \psi(u, g(u)) p(u) \, du & \text{Evaluate } \delta_{g(u)}(v)\\
		& = \textsf{RHS} & \text{Definition of pushforward}
	\end{align}
	Here we've made an assumption that $g$ doesn't take multiple values / doesn't have multiple roots; to make the proof simpler. In real life, measures which atomic will have splitting.
	
	\item \emph{We want to show} $\forall$ bounded \& measurable $\psi: \mathbb R^d \to \mathbb R$.
	\begin{align}
		\int \psi(v) p'(v) \, dv = \int \psi(v) (g_\#p)(v) \, dv
	\end{align}
	\emph{Starting proof} from the LHS
	\begin{align}
		\textsf{LHS} & = \int_{\mathbb R^d} \psi(v) \left[ \int_{\mathbb R^d} \pi(u,v) \, du \right] \, dv & \text{Existance of coupling}\\
		& = \int_{\Omega} \psi(v) \, p(u) \delta_{g(u)}(v) \, du \, dv & \text{previous property}\\
		& = \int_{\mathbb R^d} \psi(g(u)) \, p(u) \, du & \text{Evaluate } \delta_{g(u)}(v)\\
		& = \textsf{RHS} & \text{Def of pushforward}
	\end{align}
	\end{enumerate}
\end{proof}

\begin{theorem}
	If $c(u,v) = |u-v|^2$ and $p \in L^1(\mathbb R^d; \mathbb R)$ and $p,p' \in \mathcal P_2(\mathbb R^d) = \left\{ \text{has second moments} \right\}$. Then the Monge problem has a \emph{unique minimizer} $g^*$ and the Kantorovich Formulation has a \emph{unique minimizer} $\pi^* = (\textsf{id}, g^*)_\# p$.
\end{theorem}

\begin{eg}
	Let's compute the optimal transport map between two Gaussians $p_i = \mathcal N(m_i, \Sigma_i) \in \mathcal P(\mathbb R^d)$.Assuming $\Sigma_1$ is invertible.
	
	Ansatz the affine map $g(u) = Au+b$ which enables pushforward $p_2 = g_\# p_1$. So we can identify the coefficients of the affine map by computing sufficient statistics. First the mean
	\begin{align}
		m_2 = \mathbb E_{u \sim p_1} [g(u)] = A m_1 + b \implies b = m_2 - A m_1
	\end{align}
	meaning $g(u) = A(u - m_1) + m_2$. Now we can do the covariance
	\begin{align}
		\Sigma_2 & = \mathbb E_{u \sim p_1} [(g(u) - m_2)(g(u) - m_2)^T]\\
		& = A \mathbb E_{u \sim p_1}[(u - m_1)(u- m_1)^T]A^T\\
		& = A \Sigma_1 A^T
	\end{align}
	To find the solution, make the ansatz $A = \Sigma_2^{1/2} Q \Sigma_1^{-1/2}$ (this implies $Q \in \textsf{SO}(d)$). \textbf{The rest of the proof is homework}
	
	I'll denote the family of affine maps as
	\begin{align}
		g_Q(u) = m_2 + \Sigma_2^{1/2} Q \Sigma_1^{-1/2} (u- m_1)
	\end{align}
\end{eg}
\begin{align}
	W_2(p_1,p_2)^2 & = \int_{\mathbb R^d} |u- g^*(u)|^2 \, p_1(u) \, du\\
& \leq \int_{\mathbb R^d} |u - g_{\mathbb I}(u)|^2 p_1(u) \, du & \text{Suboptimal map}\\
& = \mathbb E_{u \sim p_1} |u - g_{\mathbb I}(u)|^2\\
& = |m_1 - m_2|^2 + \mathbb E_{u \sim p_1}|S(u - m_1)|^2\\
& = |m_1 - m_2|^2 + \mathbb E_{u \sim p_1} \text{Tr} \left( S(u-m_1) (u-m_1)^T S^T \right)\\
& = |m_1 - m_2|^2 + \text{Tr}(S \Sigma_1 S^T) & [\text{Trace}, \mathbb E] = 0\\
& = |m_1 - m_2|^2 + \text{Tr}\left( [\mathbb I - \Sigma_2^{1/2} \Sigma_2^{-1/2}][\Sigma_1 - \Sigma_1^{1/2} \Sigma_2^{1/2}] \right)\\
& = |m_1 - m_2|^2 + \text{Tr} \left( \Sigma_1 + \Sigma_2 \right) - 2 \text{Tr}(\Sigma_1^{1/2} \Sigma_2^{1/2})
\end{align}
\begin{remark}
	If $d=1$. Let $m_1= v, \Sigma_1 = 0$ and $m_2 = m, \Sigma_2 = \sigma^2$. Then
	\begin{align}
		W_2^2(p_1, p_2) \leq |v - m|^2 + \sigma^2
	\end{align}
	So now we can make sensible distance measures between atomic measures \& full support measures. Which was failed by the $\textsf{KL}$ and $\textsf{TV}$.
\end{remark}


\subsection{Measures over Probability Measures}
\begin{definition}
	[Kernel] Let $D \subseteq \mathbb R^d$. A kernel is a map $c : D \to  D \to \mathbb R$. It is symmetric
\end{definition}
\begin{definition}
	[Gram Matrix] IF for all $n \in \mathbb N$ and $u \in D^{\otimes^N}$ 
\end{definition}
\begin{theorem}
	[Mercer] Let $D \subseteq \mathbb R^d$ be a compact 
\end{theorem}

\subsection{Metrics over Random Probability Measures}
Sometimes it's useful to consider random variables which take values in the space of probability measures. Consider a function $\pi : \Omega \to \mathcal P(\mathbb R^d)$ for a probability space $(\Omega, \mathcal B, \mathbb P)$. For $\omega \in \Omega$, let $\pi = \pi(\omega)$; we say this is an element of the space of random probability measures.

\begin{definition}[Metric on Random Measuresj]
	Let $\pi, \pi' : \Omega \to \mathcal P(\mathbb R^d)$. Then
	\begin{align}
		d_{rand}(\pi, \pi') = \sup_{|f|_{\infty} < 1} \left | \mathbb E \left[  \left( \mathbb E_\pi[f] - \mathbb E_{\pi'}[f] \right)^2 \right] \right |^{1/2}
	\end{align}
	where $\mathbb E$ is taken over the randomness sampling probability measures.
\end{definition}
\begin{remark}
	When $\pi,\pi'$ are not random, $d_{rand}$ reduces to twice the TV distance.
\end{remark}
\begin{proof}
	Let the condition $\delta := (|f|_{\infty} \leq 1)$. Our metric on random measures is
	\begin{align}
		d_{rand}(\pi, \pi') = \sup_{\delta} \left | \mathbb E_{\pi}[f] - \mathbb E_{\pi'}[f]  \right |
	\end{align}
	Since $\pi, \pi'$ are assumed to not be random measures, we don't have the outer expectation.
	\begin{align}
		& = \sup_{\delta} \left | \int_{\mathbb R^d \times \mathbb R^d} [f(x) - f(x') ] \pi(x) \, \pi(x') \, dx \, dx'\right |
	\end{align}
\end{proof}



\section{Supervised Learning}
We now move onto supervised learning, which is a machine learning technique. \\
\\
First let's state our \textbf{data asssumptions}. In this case, I'll say we're given a dataset $\left\{ u^{(n)}, y^{(n)} \right\}_{i=1}^N$ drawn iid from $\gamma \in \mathcal P(\mathbb R^d \times \mathbb R)$ defined by the coupling: $u \sim \Upsilon$ and $\eta \sim \mathcal N(0, \mathbb I)$ (s.t. $u \perp \eta$). The data generating process is
\begin{align}
	y = \psi^\dagger(u) + \sqrt{\lambda} \eta
\end{align}
We'll call the .

The goal is to approximate $\psi^\dagger$ from data.\\
\\
How can we approximate $\psi^\dagger$? In these lectures we'll inspect a couple \textbf{approximation classes}
\begin{itemize}
	\item Gaussian processes: Tractable to inspect theoretically. 
	\item Random features
	\item Neural networks: Difficult to inspect theoretically / get guarentees.
\end{itemize}
but there are many more!\\
\\
Consider a parametric function $\psi : \mathbb R^d \times \Theta \to \mathbb R$ (s.t. $\Theta \subseteq \mathbb R^p$). Consider the population loss (assumes you have access to infinite data points)
\begin{align}
	J(\theta) & = \frac{1}{2} \mathbb E_{(U,Y) \sim \gamma} [|y - \psi(y ; \theta)|^2] & \text{Population loss}\\
	\theta^* & \in \arg \min_{\theta \in \Theta} J(\theta) & \text{Population objective}
\end{align}
Now we can inspect the loss function we have access to. Consider the empirical measure $\gamma^N = \frac{1}{N} \sum_{n=1}^N \delta_{(u^{(n)}, y^{(n)})} \in \mathcal P((\mathbb R^d \times \mathbb R)^{\otimes N})$. Note: $d_{rand} (\gamma^N, \gamma) = \frac{1}{\sqrt N}$ (monte carlo approximation). Then we have access to the empirical loss function
\begin{align}
	J^N(\theta) & = \frac{1}{2} \mathbb E_{(u,y) \sim \gamma^N} |y - \psi(u; \theta)|^2\\
	& = \frac{1}{2} \sum_{n=1}^N |y^{(n)} - \psi(u^{(n)}; \theta)|^2\\
	\theta & \in \arg \min_{\theta \in \Theta} J^N(\theta)
\end{align}

\subsection{Reproducing Kernel Hilbert Space (RKHS)}
\begin{definition}
	[Kernel] Let $D \subseteq \mathbb R^d$. A kernel $c : D \times D \to \mathbb R$ is symmetric if $c(u, u') = c(u', u)$ $\forall u, u' \in D$.
\end{definition}
Let $N \in \mathbb N$ and $U = \left\{ u^{(n)} \right\}_{n=1}^N$
\begin{definition}
	[Gram Matrix] Consider a $\tilde C \in \mathbb R^{N \times N}$ s.t. $\tilde C_{ij} = c(u^{(i)}, u^{(j)})$. If $\forall n \in \mathbb N$ and $U \in D^{\otimes^N}$, $C \succeq 0$ then kernel is resp. strictly positive definite.
\end{definition}

\begin{theorem}
	[Mercer] Let $D \subseteq \mathbb R^d$ be a compact. Let $c: D \times D \to \mathbb R$ is symmetric, and positive definite, and continuous. Then define $C : L^2(D) \to L^2 (D)$ as
	\begin{align}
		(C \varphi)(u) = \int_D c(u,u') \varphi(u') \, du'
	\end{align}
	Then $C$ has eigenpairs $(\sigma_i, \varphi_i)_{i=1}^\infty$
	\begin{itemize}
		\item Orthonormal basis: $\left\{ \varphi_i \right\}$ forms an orthonormal basis for $L^2$
		\item Finite energy: $\sum_{i=1}^\infty \sigma_i < \infty$
		\item Eigenbasis expansion: $c(u,u') = \sum_{i=1}^\infty \sigma_i \varphi_i(u) \varphi_i(u')$ 
	\end{itemize}
	This should remind you of a spectral theorem type of theorem.
\end{theorem}

\begin{definition}[RKHS]
	A Hilbert space $\mathcal H$ of function $f: D \to \mathbb R$ is a RKHS if $\exists c : D \times D \to \mathbb R$ with $c(u, \cdot) \in \mathcal H$ if $\forall u \in D$ and $\forall f \in \mathcal H$ we have the reproducing property
	\begin{align}
		f(u) = \langle f, c(u, \cdot)\rangle_{\mathcal H}
	\end{align}
\end{definition}
This connects a point-wise evaluation of a function, to a construction of a function (that is in a Hilbert Space).

\begin{eg}
	[Counter Example] Consider
	\begin{align}
		\int f(u') \delta(u-u') du' = f(u)
	\end{align}
	We'd like to say ``$\langle f, \delta_u \rangle_{L^2} = f(u)$", but $\delta_u \notin L^2$ but rather is in the space of distributions. 
\end{eg}
So instead let's define an inner product $\langle f, g \rangle_{\mathcal H} = \langle C^{-1/2} f , C^{-1/2} g\rangle_{L^2} = \langle f, C^{-1} g\rangle$, and define the $(C^{-1} c)(u, \cdot) = \delta_u$. 
\begin{eg}
	Let $C^{-1} = - \frac{d^2}{dx^2}$ by the covariance operator of a Brownian bridge then $c(\cdot, u)$ is the associated Greens function.
\end{eg}

\begin{definition}[Modified Function Space, Inner Product \& Norm]
	Under the Mercer conditions. Let $\mathcal K \subset L^2$. Define a weighted $L^2$ inner product
	\begin{align}
		\langle f, g \rangle_{\mathcal K} = \sum_{i=1}^\infty \frac{1}{\sigma_i} \langle f , \varphi_i \rangle_{L^2} \langle g, \varphi_i \rangle_{L^2}
	\end{align}
	For the norm $\|f\|_{\mathcal K} = \langle f, f\rangle_{\mathcal K}$ (induced by the inner product) to be finite, we need $\mathcal K := \left\{ f \in L^2(D) : \| f\|_{\mathcal K}^2 < \infty \right\} \subset L^2(D)$.
\end{definition}
\begin{theorem}[RKHS]
	For $f \in \mathcal K$ and $u \in D$
	\begin{align}
		\langle f , c(u, \cdot) \rangle_{\mathcal K} = f(u)
	\end{align}
\end{theorem}
\begin{proof}
	First expand $c$ using Mercer theorem
	\begin{align}
		\langle c(u, \cdot,), \varphi_i\rangle_{L^2} & = \sum_{j=1}^\infty \sigma_j \varphi_j(u) \underbrace{\langle \varphi_j, \varphi_i\rangle_{L^2}}_{=\delta_{ji}}\\
		& = \sigma_i \varphi_i(u)
	\end{align}
	Notice we have an additional $\sigma_i$, so if instead we used the modified inner product
	\begin{align}
		\langle c(u, \cdot), f\rangle_{\mathcal K} & = \sum_{i=1}^\infty \frac{1}{\sigma_i} \langle f, \varphi_i \rangle_{L^2} \langle c(u, \cdot) , \varphi_i\rangle_{L^2}\\
		& = \sum_{i=1}^\infty \langle f , \varphi_i \rangle_{L^2} \varphi_i(u)\\
		& = f(u)
	\end{align}
	In the 2nd step, I use the previous result; in the 3rd step I recall this is just the eigenbasis expansion of $f$.
\end{proof}
The terminology goes like: First pick a kernel $c$, using the Mercer theorem we get the eigensystem $(\sigma_i, \psi_i)_{i=1}^\infty$, which then can be used to define the inner product $\langle \cdot, \cdot \rangle_{\mathcal K}$ which defines a Hilbert space.

\begin{eg}
	[Squared Exponential Kernel] Let $D \subset \mathbb R^d$ be bounded \& open. Choose the kernel to be a Gaussian (in approximation theory, these are called radial basis functions)
	\begin{align}
		c(u, u') = \exp \left( - \frac{1}{2\sigma^2} |u - u'|^2 \right)
	\end{align}
	If $f \in \textsf{RKHS}$ then $f \in C^\infty(D; \mathbb R)$.
\end{eg}
\begin{theorem}
	Let $\mathcal K$ be a RKHS of a squared exponential kernel, $D \subset \mathbb R^d$ bounded \& open. Then for $\psi^\dagger \in C(D; \mathbb R)$. Then $\forall \epsilon > 0$, $\exists \psi \in \mathcal K$ s.t.
	\begin{align}
		\sup_{u \in D} |\psi^\dagger - \psi(u)| < \epsilon
	\end{align}
\end{theorem}

Now we can go back to the supervised learning problem. Consider the same problem but modified with a regularizer in function space
\begin{align}
	J^N(\psi) & = \frac{1}{2N} \sum_{n=1}^N |y^{(n)} - \psi(u^{(n)} ; \theta)|^2 + \frac{\lambda}{2N} \| \psi\|_{\mathcal K}^2\\
	\psi^* & \in \arg \inf_{\psi \in \textsf{RKHS}} J^N(\psi)
\end{align}
Somehow we've gotten worse. Now we're optimizing over infinite dimensions. However, a remarkable fact, this optimization problem is actually finite dimensional.
\begin{theorem}
	[Representer Theorem] 
	\begin{align}
		\psi^* \in \textsf{span} \left\{ c(\cdot; u^{(n)})  \right\}_{n=1}^N
	\end{align}
\end{theorem}
So now we can parameterize the space of solutions as
\begin{align}
	\psi(u; \alpha) = \sum_{n=1}^N \alpha_n c(u, u^{(n)})
\end{align}
and now just optimize on $\alpha_n$.


\subsection{Random Feature Regression} Consider the probability space $(\Omega, \mathcal F, q)$, meaning $q \in \mathcal P(\Omega)$. We'll look at random functions $\varphi: \mathcal D \times \Omega \to \mathbb R$. Note $\theta_i \in \Omega$ is the RNG seed space, so we can sample functions by $\varphi_i := \varphi(\cdot, \theta_i)$ where $\theta_i \sim q_i$. 
\begin{eg}
	[Random Fourier Feature] Let $\mathcal D = \mathbb R^d$. Consider
	\begin{align}
		\varphi(u, \theta) = \cos(\langle w , u\rangle + b), \quad \theta = \left\{ w, b \right\} 
	\end{align}
\end{eg}
A couple assumptions moving forward
\begin{enumerate}
	\item $\varphi_i \in C(\bar D ; \mathbb R)$ (for all $i \in [m]$)
	\item $\left\{ \varphi_i \right\}_{i=1}^m$ are linearly independent.
\end{enumerate}
Our objective is to find
\begin{align}
	\psi = (u ; \beta) = \frac{1}{\sqrt{m}} \sum_{i=1}^m \beta_i \varphi_i(u) \label{eqn:RandomFeature_Ansatz}
\end{align}  
where $\varphi_i$ are the randomized features, and $\beta_i$ are the parameters of the model.
\paragraph{Reproducing Kernel Hilbert Space}
Let $\mathcal K_m = \textsf{span} \left\{ \varphi_i \right\}_{i=1}^m$, and $c_m(u, u') = \frac{1}{m} \sum_{j=1}^m \varphi_j(u) \varphi_j(u')$. 
\begin{sidework}
	Looking ahead: Notice that taking the limit $m \to \infty$, and since the $\varphi_j$ are iid. You might expect $ c_m(u,u') \to c(u,u') = \mathbb E_{\theta \sim q}[\varphi(u; \theta) \varphi(u', \theta)]$.
\end{sidework}
\begin{remark}
	Under the assumptions, the equation $\ref{eqn:RandomFeature_Ansatz}$ defines an isomorphism between $\mathcal K_m \leftrightarrow \mathbb R^m$. Given $\psi \in \mathcal K_m$ we write $\beta^\psi \in \mathbb R^m$.
\end{remark}
So now we can construct a norm using the isomorphism
\begin{definition}
	Inner product $\langle f, g \rangle_{\mathcal K_m} = \langle \beta^f, \beta^g\rangle_{\mathbb R^m}$
\end{definition}
\begin{remark}
	$\|f\|_{\mathcal K_m}^2 = |\beta^f|^2$ (where $|\cdot|$ is the Euclidian norm).
\end{remark}
\begin{theorem}
	$(\mathcal K_m, \langle \cdot, \cdot \rangle_{\mathcal K_m})$ is a RKHS with kernel $c_m$ given by
	\begin{align}
		c_m(u,u') = \frac{1}{m} \sum_{j=1}^m \varphi_j(u) \varphi_j(u')
	\end{align}
\end{theorem}
Recall that $\mathcal K$ is an RKHS if (1) $c(u, \cdot) \in \textsf{RKHS}$ for all $u \in \mathcal D$, (2) $\langle f, c(u, \cdot) \rangle_{\mathcal K} = f(u)$.
\begin{proof}
	Fix $u \in \mathcal D$. By assumption (1), we can evaluate the functions pointwise
	\begin{align}
		c_m(u, \cdot) & = \frac{1}{m} \sum_{j=1}^m \varphi_j(u) \varphi_j(\cdot )\\
		& = \frac{1}{\sqrt{m}} \sum_{j=1}^m \underbrace{\frac{1}{\sqrt{m}} \varphi_j(u)}_{\beta_j^{c_m(u, \cdot)}} \varphi_j(\cdot)\\
		& = \frac{1}{\sqrt{m}} \sum_{j=1}^m \beta_j^{c_m(u, \cdot)} \varphi_j(u) \in \textsf{span}\left\{ \varphi_j \right\}_{j=1}^m
	\end{align}
	To be clean $\beta^{c_m(u, \cdot)} = \frac{1}{\sqrt{m}} (\varphi_1(u), ..., \varphi_m(u))$. This proves the first property, to prove the second property
	\begin{align}
		\langle f, c(u, \cdot) \rangle_{\mathcal K}& = \langle \beta^f, \beta^{c(u, \cdot)} \rangle_{\mathbb R^m}\\
		& = \frac{1}{\sqrt{m}} \sum_{j=1}^m \beta_j^f \varphi_j(u)\\
		& = f(u)
	\end{align}
\end{proof}

\begin{align}
	J_m(\beta) & = \frac{\lambda}{2N} | \beta|^2 + \frac{1}{2N} \sum_{n=1}^N \left| y^{(n)} - \frac{1}{\sqrt{m}} \sum_{j=1}^m \beta_j \varphi_j(u^{(n)}) \right|^2\\
	\beta^* & \in \arg \min_{\beta \in \mathbb R^m} J_m(\beta)
\end{align}

\section{Unsupervised Learning}
\paragraph{Setup} You are given a dataset $\textsf{U} = \left\{ u^{(n)} \right\}_{n=1}^N$ s.t. $u^{(n)} \overset{iid}{\sim} \gamma$. Denote $\gamma^N = \frac{1}{N} \sum_{n=1}^N \delta_{u^{(n)}}$, notice $d_{rand}(\gamma^N , \gamma) = 1/\sqrt{N}$. Your objective is to do generative modeling; find $q(\cdot ; \theta^*) \in \mathcal P(\mathbb R^d)$ such that $q(\cdot ; \theta^*) \approx \gamma$. 

\paragraph{Method} Use the loss function
\begin{align}
	F(\theta) & = \textsf{KL}(\gamma \| q(\cdot; \theta)), \quad \theta^* \in \arg \min_{\theta \in \Theta} F(\theta)
\end{align}
However, there seems to be a problem, because we only have access to $\gamma^N = \frac{1}{N} \sum_{n=1}^N \delta_{u^{(n)}}$. So the $\textsf{KL}$ is ill-defined, and additionally there's a question about generalization; how can we prevent just exactly memorizing the samples $\left\{ u^{(n)} \right\}_n$. Let's inspect the loss function closer
\begin{align}
	& = \mathbb E_{u \sim \gamma}[\log \gamma] - \mathbb E_{u \sim \gamma}[\log q(\cdot; \theta)]\\
	& = C - \mathbb E_{u \sim \gamma}[\log q(\cdot; \theta)]\\
	& = C + \bar{\textsf{LS}}[q(\cdot ; \theta), \gamma]
\end{align}
Since we're training, we only need access to derivatives of $F(\theta)$, so just use 
\begin{align}
	J(\theta) & = - \mathbb E_{\gamma} [\log q(\cdot; \theta)]\\
	J^N(\theta) & = - \mathbb E_{\gamma^N} [\log q(\cdot;  \theta)] \\
	& = - \frac{1}{N} \sum_{n=1}^N \log q(u^{(n)} ; \theta)
\end{align}
\begin{fact}
	Minimizing the KL is equivalent to maximizing the Likelihood.
\end{fact}


















\newpage
hi



\end{document}
